\chapter{STM32 Board Equivalence Search}
\label{chap:cascade-board-equivalence}

As a separate assignment during the internship, the host team provided us with a list of
competitor development boards and asked us to identify the nearest STM32 equivalent for
each one. This task was an independent market and product-positioning study. It did not
select the platforms used by the benchmark campaigns presented in the other chapters and
had no direct effect on their scenarios, measurements, or conclusions.

We recorded the outcome in the internal project artifact \emph{board equivalence.pdf}.
The proposed mappings identify the nearest STM32 board according to the characteristics
available for comparison; they are not claims that two boards are interchangeable.

\section{Meaning of board proximity}
\label{sec:cascade-equivalence-method}

For this assignment, \emph{near} means similar according to a combination of product
characteristics. We considered processing class, memory capacity, peripheral support,
development-board capabilities, and intended use. No single candidate matches every
characteristic. Architecture, memory organisation, peripheral implementation, security
features, package, price, and development tooling may remain different.

\section{Candidate mappings}
\label{sec:cascade-equivalence-mappings}

Table~\ref{tab:cascade-board-mappings} records the final candidate mappings in the internal
summary. We retain their status as proposed correspondences and do not attach measurements
to them.

\begin{table}[H]
\centering
\caption{Candidate board mappings recorded in the internal equivalence study. Official product and device sources support the subsequent identifier and lifecycle checks~\cite{nxp-imxrt1060-evk,st-stm32h735g-dk,nxp-frdm-mcxa156,st-nucleo-l552ze-q,renesas-fpb-ra0e1,st-nucleo-l031k6,renesas-fpb-ra8e1,st-nucleo-h743zi2,ti-lp-mspm0c1104,gd32e507-ds,st-nucleo-f429zi}.}
\label{tab:cascade-board-mappings}
\begin{tabular}{lll}
\toprule
Vendor & Candidate board & STM32 candidate \\
\midrule
NXP & MIMXRT1060-EVKC & STM32H735G-DK \\
NXP & FRDM-MCXA156 & NUCLEO-L552ZE-Q \\
Renesas & FPB-RA0E1 & NUCLEO-L031K6 \\
Renesas & FPB-RA8E1 & NUCLEO-H743ZI2 \\
Texas Instruments & LP-MSPM0C1104 & NUCLEO-L031K6 \\
GigaDevice & GD32E507Z-EVAL & NUCLEO-F429ZI \\
\bottomrule
\end{tabular}
\end{table}

The internal summary identifies the first NXP board as MIMXRT1060-EVKC. NXP's current
catalog page is filed under MIMXRT1060-EVKB~\cite{nxp-imxrt1060-evk}; we use it to support
the product-family check, not to assert that the two board revisions are interchangeable.

The source summary spells the first STM32 board as ``SM32H735G-DK''. The official ST
product page confirms that the intended identifier is STM32H735G-DK~\cite{st-stm32h735g-dk}.

\section{Interpretation and limitations}
\label{sec:cascade-equivalence-caveats}

Several mappings cross processor-architecture boundaries. The FPB-RA0E1 and
NUCLEO-L031K6 pair relates Cortex-M23 and Cortex-M0+ platforms; FPB-RA8E1 and
NUCLEO-H743ZI2 relate Cortex-M85 and Cortex-M7 platforms; and GD32E507Z-EVAL and
NUCLEO-F429ZI relate Cortex-M33 and Cortex-M4 platforms. Substantial memory, peripheral,
and software-ecosystem differences may also remain. The mappings should therefore be read
as the result of a nearest-match search within the STM32 portfolio, not as declarations of
technical identity or drop-in compatibility. The current ST product pages also mark the
original NUCLEO-H743ZI and NUCLEO-F429ZI boards obsolete; lifecycle status must therefore
be considered before reusing these historical proposals~\cite{st-nucleo-h743zi2,st-nucleo-f429zi}.

\section{Assignment outcome}
\label{sec:cascade-equivalence-outcome}

The assignment produced six documented competitor-to-STM32 proposals and made the
selection criteria and mismatches explicit. This deliverable was returned to the host team
as a standalone product-comparison study. The work ends at this point; the subsequent USB,
GD32, and TI benchmark campaigns use their own independently selected boards and
experimental controls.